\documentclass[11pt]{article}
\usepackage[margin=1in]{geometry}
\usepackage{xcolor}
\usepackage[breakable,listings]{tcolorbox}
\tcbuselibrary{listings}
\usepackage{hyperref}
\usepackage{enumitem}
\usepackage{listings}
\usepackage{verbatim}
\usepackage{amsmath}
\usepackage{graphicx}
\usepackage{booktabs}
\usepackage{float}

% Define colors
\definecolor{osaorange}{RGB}{220,115,38}
\definecolor{aiblue}{RGB}{30,144,255}
\definecolor{codebackground}{RGB}{248,248,248}

% Configure hyperref
\hypersetup{
    colorlinks=true,
    linkcolor=blue,
    urlcolor=blue,
    citecolor=blue
}

% Configure listings for code display
\lstset{
  basicstyle=\small\ttfamily,
  breaklines=true,
  breakatwhitespace=true,
  postbreak=\mbox{\textcolor{gray}{$\hookrightarrow$}\space},
  columns=flexible,
  keepspaces=true,
  showstringspaces=false
}

% Code environment using tcolorbox with listings
\newtcblisting{rcode}{
  colback=white,
  colframe=gray!50,
  boxrule=0.5pt,
  arc=0pt,
  left=3pt,
  right=3pt,
  top=3pt,
  bottom=3pt,
  width=\textwidth,
  breakable,
  listing only,
  listing options={
    basicstyle=\small\ttfamily,
    breaklines=true,
    breakatwhitespace=true,
    postbreak=\mbox{\textcolor{gray}{$\hookrightarrow$}\space},
    columns=flexible,
    keepspaces=true,
    showstringspaces=false
  }
}

\title{}
\author{}
\date{}

\begin{document}

% Header box
\begin{tcolorbox}[colback=osaorange,colframe=black,boxrule=2pt,arc=0pt,left=3pt,right=3pt,top=6pt,bottom=6pt,boxsep=2pt]
\centering
\textcolor{white}{\textbf{\Large H524 Introduction to Biostatistics}}\\
\textcolor{white}{\textbf{\Large Week 2 Lab: R Basics and Introduction to Probability}}\\
\textcolor{white}{\textbf{\Large Fall 2025}}
\end{tcolorbox}

\vspace{8pt}

\noindent\textbf{Format:} Online\\
\textbf{Tools Required:} R, RStudio\\
\textbf{Estimated Time:} 90 minutes

\section{Learning Objectives}

By the end of this lab, you will be able to:
\begin{enumerate}
\item Use basic R arithmetic and mathematical operations
\item Work with vectors and perform calculations on health data
\item Create and manipulate data frames for medical studies
\item Calculate simple probabilities using R
\item Build 2x2 tables for diagnostic testing
\item Use R functions for binomial and normal distributions
\item \textcolor{aiblue}{\textbf{Use AI tools to help debug R code and understand probability concepts}}
\end{enumerate}

\section{Part 1: R as a Calculator}

Let's start with the absolute basics. R can be used as a simple calculator.

\subsection{Exercise 1.1: Basic Arithmetic}

Type these commands into R and observe the results:

\begin{rcode}
# Addition
5 + 3

# Subtraction
10 - 4

# Multiplication
6 * 7

# Division
20 / 4

# Exponentiation (power)
2^3

# Order of operations (PEMDAS applies!)
(5 + 3) * 2

# Square root
sqrt(16)

# Absolute value
abs(-10)
\end{rcode}

\textbf{Try it yourself:} Calculate $\frac{120 + 140 + 135}{3}$ (the average of three blood pressure readings) using R.

\vspace{2cm}

\subsection{Exercise 1.2: Storing Values in Variables}

Instead of repeating numbers, we can store them in variables:

\begin{rcode}
# Store a single number
age <- 45
age

# Store a calculation result
bmi_numerator <- 75  # weight in kg
bmi_denominator <- 1.7^2  # height in meters, squared
bmi <- bmi_numerator / bmi_denominator
bmi

# Round to 1 decimal place
round(bmi, 1)
\end{rcode}

\textbf{Practice:} Create variables for systolic blood pressure (120) and diastolic blood pressure (80), then calculate their average (mean arterial pressure approximation).

\vspace{2cm}

\section{Part 2: Working with Vectors (Lists of Numbers)}

In health research, we rarely work with single numbers. We usually have data from multiple patients.

\subsection{Exercise 2.1: Creating Vectors}

A vector is a list of values. We create vectors using the \texttt{c()} function (combine).

\begin{rcode}
# Ages of 5 patients
ages <- c(25, 34, 45, 52, 38)
ages

# How many patients?
length(ages)

# What's the average age?
mean(ages)

# What's the oldest patient's age?
max(ages)

# What's the youngest?
min(ages)
\end{rcode}

\subsection{Exercise 2.2: More Vector Operations}

\begin{rcode}
# Systolic blood pressure readings from 7 patients
systolic_bp <- c(120, 135, 142, 118, 160, 125, 138)

# Basic statistics
mean(systolic_bp)
median(systolic_bp)
sd(systolic_bp)  # standard deviation
min(systolic_bp)
max(systolic_bp)
range(systolic_bp)  # gives min and max

# Sort the values
sort(systolic_bp)

# How many readings?
length(systolic_bp)

# Sum of all readings
sum(systolic_bp)
\end{rcode}

\textbf{Practice:} Create a vector with 6 cholesterol readings: 180, 220, 195, 240, 175, 210. Find the mean, median, and standard deviation.

\vspace{2cm}

\subsection{Exercise 2.3: Accessing Specific Values}

\begin{rcode}
# Our blood pressure data
systolic_bp <- c(120, 135, 142, 118, 160, 125, 138)

# Get the first value
systolic_bp[1]

# Get the third value
systolic_bp[3]

# Get the last value
systolic_bp[length(systolic_bp)]

# Get the first three values
systolic_bp[1:3]

# Get values in positions 2, 4, and 6
systolic_bp[c(2, 4, 6)]
\end{rcode}

\section{Part 3: Creating Data Frames}

A data frame is like a spreadsheet - it has rows and columns.

\subsection{Exercise 3.1: Building a Simple Data Frame}

\begin{rcode}
# Create a small clinical trial dataset
patient_data <- data.frame(
  patient_id = c(1, 2, 3, 4, 5),
  age = c(45, 52, 38, 61, 47),
  gender = c("M", "F", "F", "M", "M"),
  systolic_bp = c(120, 135, 142, 118, 160),
  has_diabetes = c(FALSE, TRUE, FALSE, TRUE, FALSE)
)

# View the data
patient_data

# Look at the structure
str(patient_data)

# Get summary statistics
summary(patient_data)
\end{rcode}

\subsection{Exercise 3.2: Accessing Data Frame Columns}

\begin{rcode}
# Access a column using $
patient_data$age
patient_data$systolic_bp

# Calculate mean age
mean(patient_data$age)

# Calculate mean blood pressure
mean(patient_data$systolic_bp)

# Count how many have diabetes
sum(patient_data$has_diabetes)

# What proportion have diabetes?
mean(patient_data$has_diabetes)
\end{rcode}

\section{Part 4: Introduction to Probability in R}

Now let's use R to calculate some simple probabilities!

\subsection{Exercise 4.1: Simple Probability Calculations}

\textbf{Scenario:} In a group of 100 patients, 15 have diabetes.

\begin{rcode}
# Total patients
total_patients <- 100

# Patients with diabetes
diabetes_patients <- 15

# Probability of diabetes
prob_diabetes <- diabetes_patients / total_patients
prob_diabetes

# Convert to percentage
prob_diabetes * 100

# Probability of NOT having diabetes
prob_no_diabetes <- 1 - prob_diabetes
prob_no_diabetes
\end{rcode}

\subsection{Exercise 4.2: Creating a 2x2 Table}

Let's analyze a diagnostic test for a disease.

\begin{rcode}
# Create a 2x2 table (matrix in R)
# Rows: Disease status (Yes/No)
# Columns: Test result (Positive/Negative)

test_results <- matrix(
  c(85, 5,    # Disease Yes: 85 test+, 5 test-
    10, 900), # Disease No: 10 test+, 900 test-
  nrow = 2,
  byrow = TRUE,
  dimnames = list(
    Disease = c("Yes", "No"),
    Test = c("Positive", "Negative")
  )
)

# View the table
test_results

# Add row and column totals
addmargins(test_results)
\end{rcode}

\subsection{Exercise 4.3: Calculating Sensitivity and Specificity}

\begin{rcode}
# From our table above:
# True Positives (TP): disease Yes, test Positive
TP <- test_results[1, 1]

# False Negatives (FN): disease Yes, test Negative
FN <- test_results[1, 2]

# False Positives (FP): disease No, test Positive
FP <- test_results[2, 1]

# True Negatives (TN): disease No, test Negative
TN <- test_results[2, 2]

# Sensitivity: Pr(Test+ | Disease+)
sensitivity <- TP / (TP + FN)
cat("Sensitivity:", round(sensitivity * 100, 1), "%\n")

# Specificity: Pr(Test- | Disease-)
specificity <- TN / (TN + FP)
cat("Specificity:", round(specificity * 100, 1), "%\n")

# Positive Predictive Value: Pr(Disease+ | Test+)
PPV <- TP / (TP + FP)
cat("PPV:", round(PPV * 100, 1), "%\n")

# Negative Predictive Value: Pr(Disease- | Test-)
NPV <- TN / (TN + FN)
cat("NPV:", round(NPV * 100, 1), "%\n")
\end{rcode}

\section{Part 5: Probability Distributions in R}

R has built-in functions for working with probability distributions.

\subsection{Exercise 5.1: Binomial Distribution}

\textbf{Scenario:} A treatment has a 70\% success rate. We treat 10 patients.

\begin{rcode}
# Probability of exactly 7 successes
dbinom(7, size = 10, prob = 0.7)

# Probability of 7 OR MORE successes
# Method 1: Add them up
dbinom(7, 10, 0.7) + dbinom(8, 10, 0.7) +
  dbinom(9, 10, 0.7) + dbinom(10, 10, 0.7)

# Method 2: Use pbinom (much easier!)
pbinom(6, size = 10, prob = 0.7, lower.tail = FALSE)

# Expected number of successes
expected <- 10 * 0.7
expected

# Visualize the distribution
x <- 0:10  # Possible number of successes
probs <- dbinom(x, size = 10, prob = 0.7)

barplot(probs,
        names.arg = x,
        xlab = "Number of Successes",
        ylab = "Probability",
        main = "Binomial Distribution (n=10, p=0.7)",
        col = "lightblue")
\end{rcode}

\subsection{Exercise 5.2: Normal Distribution}

\textbf{Scenario:} Adult cholesterol levels are normally distributed with mean 200 mg/dL and SD 40 mg/dL.

\begin{rcode}
# What proportion have cholesterol < 240 mg/dL?
pnorm(240, mean = 200, sd = 40)

# What proportion have cholesterol > 240 mg/dL?
pnorm(240, mean = 200, sd = 40, lower.tail = FALSE)

# What proportion have cholesterol between 180 and 220?
pnorm(220, 200, 40) - pnorm(180, 200, 40)

# What cholesterol level is at the 90th percentile?
qnorm(0.90, mean = 200, sd = 40)

# Visualize the distribution
x_values <- seq(100, 300, by = 1)
y_values <- dnorm(x_values, mean = 200, sd = 40)

plot(x_values, y_values,
     type = "l",
     xlab = "Cholesterol (mg/dL)",
     ylab = "Density",
     main = "Normal Distribution of Cholesterol Levels",
     col = "blue",
     lwd = 2)

# Add vertical line at the mean
abline(v = 200, col = "red", lty = 2, lwd = 2)
\end{rcode}

\section{Practice Problems}

\subsection{Problem 1: Blood Pressure Analysis}

Seven patients have the following systolic blood pressure readings: 115, 128, 142, 135, 120, 155, 138.

\begin{enumerate}[label=\alph*)]
\item Create a vector with these values
\item Calculate the mean, median, and standard deviation
\item How many patients have high blood pressure ($\geq$ 140 mmHg)?
\item What proportion have high blood pressure?
\end{enumerate}

\vspace{3cm}

\subsection{Problem 2: Diagnostic Test}

A rapid flu test was given to 500 patients. The results were:
\begin{itemize}
\item 40 patients had the flu and tested positive
\item 5 patients had the flu and tested negative
\item 30 patients didn't have the flu but tested positive
\item 425 patients didn't have the flu and tested negative
\end{itemize}

\begin{enumerate}[label=\alph*)]
\item Create a 2x2 matrix with these results
\item Calculate the sensitivity
\item Calculate the specificity
\item Calculate the PPV (positive predictive value)
\end{enumerate}

\vspace{3cm}

\subsection{Problem 3: Treatment Success}

A new medication has an 80\% success rate. A clinic treats 15 patients with this medication.

\begin{enumerate}[label=\alph*)]
\item What is the probability that exactly 12 patients respond successfully?
\item What is the probability that at least 12 patients respond successfully?
\item What is the expected number of successful responses?
\item Create a barplot showing the probability distribution
\end{enumerate}

\vspace{3cm}

\subsection{Problem 4: BMI Distribution}

Adult BMI is normally distributed with mean 26.5 and standard deviation 5.

\begin{enumerate}[label=\alph*)]
\item What proportion of adults have BMI $<$ 25 (normal weight)?
\item What proportion have BMI $\geq$ 30 (obese)?
\item What BMI value represents the 75th percentile?
\item Plot the normal distribution curve
\end{enumerate}

\vspace{3cm}

\section{Saving Your Work}

\textbf{IMPORTANT:} Save your R script file!

\begin{rcode}
# At the top of your R script, add comments with your name and date:
# Name: Your Name
# Date: Today's Date
# Lab: Week 2 - R Basics and Probability

# Save your workspace (all objects you created)
save.image("week2_lab.RData")

# To load it later:
# load("week2_lab.RData")
\end{rcode}

\section{\textcolor{aiblue}{AI-Enhanced Learning Tips}}

\begin{tcolorbox}[colback=aiblue!10,colframe=aiblue,title=Using AI Tools for R Learning]

\textbf{Good prompts for AI assistance:}
\begin{itemize}
\item ``Explain what this R error means: [paste error message]''
\item ``How do I create a vector in R with values 1 through 10?''
\item ``Show me how to calculate sensitivity from a 2x2 table in R''
\item ``What's the difference between dbinom and pbinom in R?''
\end{itemize}

\textbf{Remember:}
\begin{itemize}
\item Always test AI-generated code before using it
\item Make sure you understand what the code does
\item AI is a learning tool, not a replacement for understanding
\end{itemize}
\end{tcolorbox}

\section{Additional Resources}

\begin{itemize}
\item R Documentation: Type \texttt{?function\_name} in R (e.g., \texttt{?mean})
\item Quick-R: \url{https://www.statmethods.net/}
\item R for Data Science (free online): \url{https://r4ds.had.co.nz/}
\item Pagano \& Gauvreau textbook: Chapter 6 (Probability)
\end{itemize}

\end{document}
